% =====================================================================
%  Panduan Pengguna Apoteku
%  Kompilasi:  tectonic panduan-pengguna.tex
% =====================================================================
\documentclass[11pt,a4paper]{article}

\usepackage[T1]{fontenc}
\usepackage[utf8]{inputenc}
\usepackage{lmodern}
\usepackage[english,bahasa]{babel}
\usepackage{graphicx}
\usepackage[font=small,labelfont=bf,skip=4pt]{caption}
\usepackage[a4paper,margin=2.4cm,top=2.6cm,bottom=2.6cm]{geometry}
\usepackage{xcolor}
\usepackage{booktabs}
\usepackage{array}
\usepackage{tabularx}
\usepackage{longtable}
\usepackage{enumitem}
\usepackage{titlesec}
\usepackage{fancyhdr}
\usepackage[colorlinks=true,linkcolor=primary,urlcolor=primary,citecolor=primary]{hyperref}
\usepackage[most]{tcolorbox}

% ---------------------------------------------------------------- warna
\definecolor{primary}{HTML}{0F766E}   % teal
\definecolor{accent}{HTML}{B45309}    % amber tua
\definecolor{ink}{HTML}{1F2937}
\definecolor{softgray}{HTML}{F3F4F6}
\definecolor{rulegray}{HTML}{D1D5DB}

% ------------------------------------------------------------- perintah
\newcommand{\app}{Apoteku}
\newcommand{\btn}[1]{\textcolor{primary}{\textbf{#1}}}
\newcommand{\field}[1]{\textit{#1}}
\newcommand{\code}[1]{\texttt{\small #1}}

% Sisipkan screenshot (docs/images) selebar area teks, dengan keterangan bernomor.
\newcommand{\screenshot}[2]{%
  \begin{figure}[htbp]
    \centering
    \includegraphics[width=\linewidth]{images/#2}
    \caption{#1}
  \end{figure}%
}

% ---------------------------------------------------------------- tata letak
\titleformat{\section}{\Large\bfseries\color{primary}}{\thesection}{0.6em}{}
\titleformat{\subsection}{\large\bfseries\color{ink}}{\thesubsection}{0.6em}{}
\titleformat{\subsubsection}{\normalsize\bfseries\color{ink}}{\thesubsubsection}{0.6em}{}
\titlespacing*{\section}{0pt}{1.6em}{0.6em}

\setlength{\parindent}{0pt}
\setlength{\parskip}{0.5em}

\pagestyle{fancy}
\fancyhf{}
\fancyhead[L]{\small\textcolor{ink}{Panduan Pengguna \app}}
\fancyhead[R]{\small\textcolor{ink}{\leftmark}}
\fancyfoot[C]{\small\textcolor{ink}{\thepage}}
\renewcommand{\headrulewidth}{0.4pt}
\renewcommand{\footrulewidth}{0pt}
\renewcommand{\sectionmark}[1]{\markboth{#1}{}}

% ---------------------------------------------------------------- kotak
\newtcolorbox{itembox}[1]{%
  colback=primary!4, colframe=primary!55, boxrule=0.5pt, arc=1.4mm,
  left=2.4mm, right=2.4mm, top=1.8mm, bottom=1.8mm, title={\textbf{#1}}, fonttitle=\color{white},
  coltitle=white, colbacktitle=primary!75, breakable}

\newtcolorbox{tipbox}[1]{%
  colback=accent!5, colframe=accent!70!black, boxrule=0.5pt, arc=1.4mm,
  left=2.4mm, right=2.4mm, top=1.8mm, bottom=1.8mm, title={\textbf{Tips: #1}}, fonttitle=\color{white},
  coltitle=white, colbacktitle=accent!85, breakable}

\newtcolorbox{warnbox}[1]{%
  colback=red!4, colframe=red!60!black, boxrule=0.5pt, arc=1.4mm,
  left=2.4mm, right=2.4mm, top=1.8mm, bottom=1.8mm, title={\textbf{Perhatian: #1}}, fonttitle=\color{white},
  coltitle=white, colbacktitle=red!70!black, breakable}

% ---------------------------------------------------------------- tabel
\newcolumntype{L}[1]{>{\raggedright\arraybackslash}p{#1}}
\newcolumntype{Y}{>{\raggedright\arraybackslash}X}
\renewcommand{\arraystretch}{1.25}

% Beri ruang lebih banyak untuk gambar agar tidak menumpuk di akhir dokumen.
\renewcommand{\topfraction}{0.92}
\renewcommand{\bottomfraction}{0.85}
\renewcommand{\textfraction}{0.08}
\renewcommand{\floatpagefraction}{0.72}
\setcounter{topnumber}{3}
\setcounter{bottomnumber}{2}
\setcounter{totalnumber}{4}

% =====================================================================
\begin{document}

% ---------------------------------------------------------- halaman judul
\begin{titlepage}
  \centering
  \vspace*{2.2cm}
  {\color{primary}\rule{\textwidth}{2pt}}\\[0.5cm]
  {\Huge\bfseries Panduan Pengguna\\[0.25cm] \app}\\[0.5cm]
  {\large Sistem Manajemen Apotek --- Aplikasi Web}\\[0.35cm]
  {\color{primary}\rule{\textwidth}{2pt}}\\[2.2cm]

  \begin{tabular}{@{}r@{\hspace{1cm}}l@{}}
    \textbf{Versi dokumen} & 1.0 \\
    \textbf{Versi aplikasi} & 0.1.0 (MVP apotek tunggal) \\
    \textbf{Tanggal}        & 16 Oktober 2026 \\
    \textbf{Untuk}          & Pemilik, apoteker, kasir, dan staf apotek \\
  \end{tabular}

  \vfill
  \begin{minipage}{0.86\textwidth}
    \small\itshape
    Dokumen ini menjelaskan cara memakai aplikasi \app{} sehari-hari: mulai dari
    masuk, mengelola produk dan persediaan, membuka shift, melayani penjualan di
    kasir, mencetak struk, sampai menutup shift dan memeriksa rekonsiliasi kas.
  \end{minipage}
  \vspace{1.2cm}
\end{titlepage}

% ---------------------------------------------------------------- daftar isi
{\hypersetup{linkcolor=ink}
\tableofcontents
\newpage}

% =====================================================================
\section{Tentang panduan ini}

Panduan ini ditujukan untuk \textbf{pengguna} aplikasi \app{} --- bukan untuk
pengembang. Seluruh langkah ditulis mengikuti tampilan aplikasi yang
sesungguhnya.

\subsection{Konvensi penulisan}

\begin{tabularx}{\textwidth}{@{}L{3.5cm}Y@{}}
  \toprule
  \textbf{Contoh penulisan} & \textbf{Artinya} \\
  \midrule
  \btn{Buka shift}   & Tombol yang harus diklik. \\
  \field{Dibayar}    & Nama kolom atau kotak isian pada formulir. \\
  \code{INV-20261016-ABC123} & Contoh kode/format data dari sistem. \\
  \midrule
  \end{tabularx}

\vspace{0.3cm}

\begin{tipbox}{}
Tips berisi cara yang lebih cepat/efisien.
\end{tipbox}

\begin{warnbox}{}
Perhatian menandai langkah yang, jika salah, dapat menyebabkan selisih kas atau
stok.
\end{warnbox}

% =====================================================================
\section{Mengenal Apoteku}

\subsection{Apa itu Apoteku}
\app{} adalah aplikasi web untuk mengelola operasional satu apotek: katalog obat,
persediaan berbasis batch dan tanggal kedaluwarsa, pembelian dari supplier,
penjualan di kasir, resep dokter, shift kasir, sampai laporan ringkas harian.
Aplikasi dapat dibuka dari komputer, tablet, maupun ponsel karena tampilannya
responsif.

\subsection{Dua hal yang paling sering membingungkan}

\textbf{1.\ Stok tidak disimpan pada produk, melainkan pada \emph{batch}.}
Setiap penerimaan barang membuat batch baru berisi jumlah dan tanggal
kedaluwarsa. Karena itu, satu produk bisa punya beberapa batch sekaligus.

\textbf{2.\ Kasir wajib membuka \emph{shift} sebelum bertransaksi.}
Shift adalah sesi kerja kasir, tempat uang awal dan uang akhir laci dicatat.
Tanpa shift terbuka, tombol pembayaran di halaman Kasir tidak akan bekerja.

\subsection{Peran pengguna}

\begin{tabularx}{\textwidth}{@{}L{2.6cm}Y@{}}
  \toprule
  \textbf{Peran} & \textbf{Singkatnya} \\
  \midrule
  \textbf{owner}      & Pemilik. Akses paling luas, termasuk pengaturan apotek dan log audit. \\
  \textbf{admin}      & Administrator. Setara owner untuk operasional harian, termasuk log audit. \\
  \textbf{pharmacist} & Apoteker. Mengelola katalog obat, pembelian, penyesuaian stok, dan resep. \\
  \textbf{cashier}    & Kasir. Melayani penjualan, membuka/menutup shift, dan mendaftarkan pelanggan. \\
  \bottomrule
\end{tabularx}

\begin{warnbox}{hak akses ditegakkan di server}
Tombol dan menu tampil untuk semua peran. Bila kasir menekan tombol yang
khusus owner/apoteker, sistem akan menolak dengan pesan \textbf{``Tidak punya
akses''}. Rincian lengkap per peran ada di Bagian~\ref{sec:rbac}.
\end{warnbox}

\subsection{Peta halaman}

\begin{tabularx}{\textwidth}{@{}L{3.4cm}Y@{}}
  \toprule
  \textbf{Menu} & \textbf{Kegunaan} \\
  \midrule
  \textbf{Dashboard}    & Ringkasan penjualan, laba, nilai persediaan, produk terlaris hari ini. \\
  \textbf{Kasir}        & Melayani penjualan (POS) dan mencetak struk. \\
  \textbf{Transaksi}    & Riwayat penjualan, cetak ulang struk. \\
  \textbf{Produk}       & Katalog obat/alkes/BHP, harga, satuan bertingkat. \\
  \textbf{Persediaan}   & Stok per produk, daftar batch, kedaluwarsa, peringatan. \\
  \textbf{Pembelian}    & Purchase order (PO) ke supplier dan penerimaan barang. \\
  \textbf{Resep}        & Pencatatan resep dokter. \\
  \textbf{Shift Kasir}  & Buka/tutup shift dan rekonsiliasi uang laci. \\
  \textbf{Pengaturan}   & Profil apotek, apoteker penanggung jawab, catatan kaki struk. \\
  \bottomrule
\end{tabularx}

% =====================================================================
\section{Memulai}

\subsection{Membuka aplikasi}
Buka peramban (Chrome, Safari, Edge) dan akses alamat aplikasi yang diberikan
oleh tim IT Anda, misalnya \code{http://localhost:8080}. Halaman masuk akan
tampil.

\subsection{Masuk (login)}

\begin{enumerate}[leftmargin=1.4em,itemsep=0.35em]
  \item Isi \field{Email}.
  \item Isi \field{Password}.
  \item Klik \btn{Masuk}. Bila sedang diproses, tombol berubah menjadi
        \btn{Memproses...}.
\end{enumerate}

\begin{itembox}{Akun contoh (data awal)}
\begin{tabularx}{\textwidth}{@{}Y L{3.2cm}Y@{}}
  \toprule
  \textbf{Email} & \textbf{Peran} & \textbf{Password} \\
  \midrule
  \code{owner@apoteku.test}    & owner      & \code{password123} \\
  \code{kasir@apoteku.test}    & cashier    & \code{password123} \\
  \code{apoteker@apoteku.test} & pharmacist & \code{password123} \\
  \bottomrule
\end{tabularx}
\end{itembox}

\begin{warnbox}{ganti password akun contoh}
Akun di atas hanya untuk masa uji coba. Sebelum dipakai sungguhan, pastikan
password sudah diganti oleh tim IT.
\end{warnbox}

\screenshot{Halaman masuk (\emph{login}) aplikasi Apoteku.}{01-halaman-masuk}

\subsection{Tata letak antarmuka}

Setelah masuk, layar terdiri dari dua bagian:

\begin{tabularx}{\textwidth}{@{}L{4.2cm}Y@{}}
  \toprule
  \textbf{Bagian} & \textbf{Keterangan} \\
  \midrule
  \textbf{Sidebar kiri} & Menu navigasi: Dashboard, Kasir, Transaksi, Produk, Persediaan, Pembelian, Resep, Shift Kasir, Pengaturan. Di bawahnya tertera nama dan peran Anda, serta tombol \btn{Keluar}. \\
  \textbf{Area kerja}   & Isi halaman yang sedang dibuka. \\
  \bottomrule
\end{tabularx}

Pada layar kecil (ponsel), sidebar berubah menjadi deretan menu di bagian atas,
dan tombol \btn{Keluar} pindah ke kanan atas.

\subsection{Keluar}
Klik \btn{Keluar} di sidebar (atau di kanan atas pada ponsel). Anda akan kembali
ke halaman masuk.

\begin{tipbox}{kenapa tiba-tiba kembali ke halaman masuk?}
Bila sesi Anda berakhir atau token kedaluwarsa, aplikasi otomatis
mengembalikan ke halaman masuk. Cukup masuk kembali.
\end{tipbox}

% =====================================================================
\section{Alur kerja harian (ringkas)}

Urutan berikut adalah cara pakai yang disarankan untuk satu hari operasional.

\begin{enumerate}[leftmargin=1.6em,itemsep=0.4em]
  \item \textbf{Pagi --- periksa Dashboard.} Lihat kartu ``Perlu perhatian'':
        berapa produk berstok rendah dan berapa batch mendekati kedaluwarsa.
  \item \textbf{Periksa Persediaan.} Buka tab \btn{Peringatan}, catat item yang
        harus segera dipesan atau dipisahkan.
  \item \textbf{Buat Pembelian} untuk barang yang menipis/hampir kedaluwarsa
        (Bagian~\ref{sec:pembelian}).
  \item \textbf{Terima barang} saat kiriman datang, agar stok dan batch baru
        tercatat.
  \item \textbf{Buka shift} di halaman Shift Kasir, isi uang awal laci
        (Bagian~\ref{sec:shift}).
  \item \textbf{Layani penjualan} di halaman Kasir, cetak struk untuk pembeli
        (Bagian~\ref{sec:kasir}).
  \item \textbf{Catat resep} yang masuk dari pelanggan (Bagian~\ref{sec:resep}).
  \item \textbf{Tutup shift} di akhir giliran kerja: masukkan uang fisik yang
        dihitung, lalu simpan (Bagian~\ref{sec:shift}).
  \item \textbf{Periksa Transaksi} untuk mencetak ulang struk
        (Bagian~\ref{sec:transaksi}).
\end{enumerate}

% =====================================================================
\section{Dashboard}
\textbf{Judul halaman:} \emph{Dashboard --- ``Ringkasan operasional apotek hari ini''}

Dashboard adalah ringkasan cepat kondisi apotek hari ini. Tidak ada tombol
aksi di halaman ini.

\screenshot{Dashboard: kartu statistik dan tabel ringkasan hari ini.}{dashboard}

\subsection{Kartu statistik}

\begin{tabularx}{\textwidth}{@{}L{4.4cm}Y@{}}
  \toprule
  \textbf{Kartu} & \textbf{Isinya} \\
  \midrule
  \textbf{Penjualan hari ini} & Total omzet hari ini, dengan keterangan jumlah transaksi. \\
  \textbf{Laba hari ini}      & Perkiraan laba (omzet dikurangi harga pokok barang terjual). \\
  \textbf{Nilai persediaan}   & Nilai stok saat ini, dengan keterangan jumlah produk aktif. \\
  \textbf{Perlu perhatian}    & Dua angka: jumlah produk berstok rendah dan jumlah batch mendekati kedaluwarsa. \\
  \bottomrule
\end{tabularx}

\subsection{Dua tabel}
\begin{itemize}[leftmargin=1.4em,itemsep=0.3em]
  \item \textbf{Produk terlaris hari ini} --- kolom \field{Produk}, \field{Qty},
        \field{Omzet}.
  \item \textbf{Transaksi terbaru} --- kolom \field{Kode}, \field{Waktu},
        \field{Total}.
\end{itemize}

\begin{tipbox}{mulai dari sini}
Kartu ``Perlu perhatian'' pada Dashboard adalah pintu masuk tercepat untuk tahu
apa yang harus dikerjakan hari ini.
\end{tipbox}

% =====================================================================
\section{Produk \& Obat}
\textbf{Judul halaman:} \emph{Produk \& Obat --- ``Master data obat, alkes, dan BHP''}

Halaman ini adalah katalog barang. Setiap penjualan, pembelian, dan resep
mengacu ke produk di sini, jadi datanya harus benar lebih dulu.

\subsection{Mencari produk}
Ketik pada kotak pencarian \field{``Cari nama / SKU / barcode...''}. Tabel
menampilkan kolom \field{Nama}, \field{SKU}, \field{Golongan}, \field{Stok}, dan
\field{Harga jual}. Bila stok sama dengan atau di bawah stok minimum, angkanya
berwarna kuning sebagai peringatan.

\screenshot{Daftar produk beserta stok dan harga jualnya.}{produk}

\subsection{Menambah produk baru}

\begin{enumerate}[leftmargin=1.4em,itemsep=0.35em]
  \item Klik \btn{Tambah produk}.
  \item Isi formulir berikut.
  \item Klik \btn{Simpan produk}. Untuk membatalkan, klik \btn{Batal} atau
        \btn{Tutup}.
\end{enumerate}

\begin{tabularx}{\textwidth}{@{}L{4.6cm}L{1.6cm}Y@{}}
  \toprule
  \textbf{Kolom isian} & \textbf{Wajib} & \textbf{Keterangan} \\
  \midrule
  \field{Nama produk} & Ya & Nama yang tampil di kasir dan struk, mis. ``Paracetamol 500mg''. \\
  \field{SKU} & Ya & Kode unik internal. Contoh format: \code{OBT-0001}. Tidak boleh sama dengan produk lain. \\
  \field{Barcode} & Tidak & Bila ada, dapat dipakai untuk pemindaian di kasir. \\
  \field{Kategori} & Tidak & Mis. Analgesik, Antibiotik, Vitamin \& Suplemen. \\
  \field{Bentuk sediaan} & Ya & tablet, kapsul, sirup, salep, krim, gel, tetes, injeksi, suppositoria, serbuk, lainnya. \\
  \field{Golongan} & Ya & obat bebas, obat bebas terbatas, obat keras, narkotika, psikotropika, alkes, bhp. \\
  \field{Satuan dasar} & Ya & Satuan terkecil untuk perhitungan stok, mis. \code{tablet}. \\
  \field{Harga jual / satuan dasar} & Ya & Harga untuk satu satuan dasar. \\
  \field{Harga pokok} & Ya & Harga beli, dipakai untuk menghitung laba. \\
  \field{Stok minimum} & Tidak & Batas peringatan stok rendah. \\
  \field{Wajib resep} & Tidak & Bila dicentang, produk ditandai \textbf{Resep} di kasir. \\
  \bottomrule
\end{tabularx}

\screenshot{Formulir tambah produk, termasuk blok ``Satuan \& konversi''.}{produk-form}

\subsection{Satuan bertingkat (mis. strip dan box)}

Pada blok \field{Satuan \& konversi}, klik \btn{Satuan} untuk menambah baris.
Isi \field{Nama} (mis. \code{strip}), \field{Konversi}, dan \field{Harga jual}.

\begin{itembox}{arti kolom Konversi}
\field{Konversi} adalah \textbf{jumlah satuan dasar} di dalam satu satuan ini.
Contoh: 1 strip berisi 10 tablet $\rightarrow$ konversi = 10. Bila harga satu
strip Rp6.500, maka saat menjual 1 strip, stok berkurang 10 tablet dan pembeli
dibayar Rp6.500.
\end{itembox}

Contoh pengisian untuk Paracetamol:

\begin{tabularx}{\textwidth}{@{}L{3cm}L{2.4cm}L{3.6cm}Y@{}}
  \toprule
  \textbf{Nama} & \textbf{Konversi} & \textbf{Harga jual} & \textbf{Artinya} \\
  \midrule
  (satuan dasar) & 1   & Rp700     & 1 tablet \\
  \code{strip}   & 10  & Rp6.500   & 1 strip = 10 tablet \\
  \code{box}     & 100 & Rp60.000  & 1 box = 100 tablet \\
  \bottomrule
\end{tabularx}

\begin{warnbox}{satu satuan dasar wajib ada}
Pastikan satuan dasar (konversi = 1) selalu ada. Satuan dasar adalah acuan
perhitungan stok di seluruh aplikasi.
\end{warnbox}

\begin{tipbox}{golongan obat}
Golongan ``obat keras'' tampil sebagai label oranye di daftar produk, agar mudah
dikenali saat melayani pembeli.
\end{tipbox}

\subsection{Catatan}
Saat ini perubahan (edit) dan penghapusan produk belum tersedia lewat tombol di
halaman ini; fitur tersebut masih melalui tim IT/API. Produk yang sudah dipakai
pada transaksi tidak dapat dihapus karena akan merusak riwayat.

% =====================================================================
\section{Persediaan}
\textbf{Judul halaman:} \emph{Persediaan --- ``Stok, batch, tanggal kedaluwarsa, dan peringatan''}

Halaman ini bersifat \textbf{lihat saja} (tidak mengubah data) dan terbagi
menjadi tiga tab. Kotak pencarian \field{``Cari produk / batch...''} tersedia
pada tab Stok dan Batch.

\subsection{Tab \btn{Stok Produk}}
Ringkasan stok tiap produk. Kolom: \field{Produk}, \field{Stok},
\field{Minimum}, \field{HPP}, \field{Status}.

\begin{tabularx}{\textwidth}{@{}L{3.2cm}Y@{}}
  \toprule
  \textbf{Status} & \textbf{Artinya} \\
  \midrule
  \textbf{Aman}       & Stok di atas batas minimum. \\
  \textbf{Stok rendah} & Stok sudah mencapai atau di bawah batas minimum --- segera rencanakan pembelian. \\
  \textbf{Habis}      & Stok nol. Produk tidak bisa dijual sebelum ada penerimaan barang. \\
  \bottomrule
\end{tabularx}

\screenshot{Tab \emph{Stok Produk}: status Aman, Stok rendah, atau Habis.}{persediaan-stok}

\subsection{Tab \btn{Batch \& Expiry}}
Rincian per batch. Kolom: \field{Produk}, \field{No. Batch},
\field{Kedaluwarsa}, \field{Sisa}, \field{Estimasi}. Tersedia filter:
\field{Semua batch}, \field{Masih ada stok}, \field{Mendekati ED},
\field{Sudah ED}, \field{Stok habis}.

\begin{tabularx}{\textwidth}{@{}L{3.6cm}Y@{}}
  \toprule
  \textbf{Tanda kedaluwarsa} & \textbf{Artinya} \\
  \midrule
  \textbf{ED lewat}          & Batch sudah kedaluwarsa. Jangan dijual; pisahkan dari rak. \\
  \textbf{N hari lagi} (kuning) & Tinggal $\le$ 90 hari. Prioritaskan batch ini agar tidak terbuang. \\
  \textbf{N hari lagi} (hijau)  & Masih lebih dari 90 hari. \\
  \textbf{--}                & Batch tidak memiliki tanggal kedaluwarsa. \\
  \bottomrule
\end{tabularx}

\screenshot{Tab \emph{Batch \& Expiry}: sisa tiap batch dan estimasi kedaluwarsa.}{persediaan-batch}

\subsection{Tab \btn{Peringatan}}
Dua daftar siap tindak:
\begin{itemize}[leftmargin=1.4em,itemsep=0.3em]
  \item \textbf{Stok di bawah minimum} --- kolom \field{Produk}, \field{Stok},
        \field{Min}.
  \item \textbf{Mendekati / melewati kedaluwarsa} --- kolom \field{Produk},
        \field{Batch}, \field{ED}, \field{Sisa}.
\end{itemize}

\screenshot{Tab \emph{Peringatan}: stok menipis dan batch mendekati kedaluwarsa.}{persediaan-peringatan}

\begin{itembox}{FEFO: batch mana yang dijual lebih dulu?}
Aplikasi otomatis mengambil stok dari \textbf{batch yang paling cepat
kedaluwarsa} (First Expired First Out). Anda tidak perlu memilih batch secara
manual saat menjual. Karena itu, pastikan tanggal kedaluwarsa diisi dengan benar
saat menerima barang.
\end{itembox}

\begin{warnbox}{batas 90 hari bersifat tetap}
Peringatan ``mendekati kedaluwarsa'' muncul bila sisa waktu $\le$ 90 hari.
Ini pengaturan bawaan sistem.
\end{warnbox}

\subsection{Catatan}
Penyesuaian stok (opname, barang rusak, barang kedaluwarsa) belum memiliki tombol
di halaman ini; saat ini dilakukan melalui tim IT/API. Bila Anda menemukan
selisih stok, catat dan laporkan agar dicatat sebagai penyesuaian.

% =====================================================================
\section{Pembelian}
\label{sec:pembelian}
\textbf{Judul halaman:} \emph{Pembelian --- ``Purchase order dan penerimaan barang dari PBF/supplier''}

\screenshot{Daftar purchase order beserta statusnya.}{pembelian}

\subsection{Membuat purchase order (PO)}

\begin{enumerate}[leftmargin=1.4em,itemsep=0.35em]
  \item Klik \btn{Buat PO}.
  \item Pilih \field{Supplier}. Bila belum ada, minta tim IT menambahkannya.
  \item Isi \field{Tanggal order}.
  \item Pada blok \field{Item pembelian}, klik \btn{Item} lalu lengkapi tiap
        baris: produk, \field{No. batch}, tanggal kedaluwarsa, \field{Qty}, dan
        \field{Harga}.
  \item Klik \btn{Simpan PO}. Status PO menjadi \textbf{ordered}.
\end{enumerate}

\begin{warnbox}{isi nomor batch dan tanggal kedaluwarsa}
Kedua kolom ini menentukan cara sistem memilih stok saat menjual (FEFO) dan
kapan peringatan kedaluwarsa muncul. Jangan dikosongkan bila barang memiliki
kedaluwarsa.
\end{warnbox}

\screenshot{Formulir purchase order: supplier, tanggal, dan item pembelian.}{pembelian-po}

\subsection{Menerima barang}
Saat kiriman datang, cari PO berstatus \textbf{ordered} pada tabel, lalu klik
\btn{Terima}. Sistem akan membuat batch baru dan menambah stok secara otomatis.
Status PO berubah menjadi \textbf{received}.

\begin{tabularx}{\textwidth}{@{}L{2.8cm}Y@{}}
  \toprule
  \textbf{Status PO} & \textbf{Artinya} \\
  \midrule
  \textbf{draft}     & Belum dikirim ke supplier. \\
  \textbf{ordered}   & Sudah dipesan; tombol \btn{Terima} tersedia. \\
  \textbf{received}  & Barang sudah diterima; stok dan batch sudah bertambah. \\
  \textbf{cancelled} & Dibatalkan. \\
  \bottomrule
\end{tabularx}

\begin{tipbox}{urutan yang benar}
Buat PO $\rightarrow$ terima barang $\rightarrow$ baru dijual. Bila barang dijual
tanpa diterima lebih dulu, stok tidak akan bertambah dan kasir akan menerima
pesan ``Stok tidak mencukupi''.
\end{tipbox}

% =====================================================================
\section{Shift Kasir}
\label{sec:shift}
\textbf{Judul halaman:} \emph{Shift Kasir --- ``Buka/tutup shift dan rekonsiliasi uang tunai laci''}

Shift adalah sesi kerja kasir. Semua penjualan yang Anda lakukan tercatat pada
shift yang sedang terbuka.

\subsection{Membuka shift}

\begin{enumerate}[leftmargin=1.4em,itemsep=0.35em]
  \item Buka menu \textbf{Shift Kasir}.
  \item Pada kotak \field{Uang awal di laci}, isi jumlah uang tunai yang Anda
        taruh di laci saat mulai bekerja (tanpa titik atau koma).
  \item Isi \field{Catatan} bila perlu (opsional).
  \item Klik \btn{Buka shift}. Status berubah menjadi \textbf{Terbuka}.
\end{enumerate}

\begin{warnbox}{hitung uang awal dengan benar}
Uang awal adalah titik nol perhitungan kas. Bila salah memasukkan angka, selisih
kas di akhir shift akan terlihat meski uang fisiknya benar.
\end{warnbox}

\screenshot{Formulir buka shift: mengisi uang awal di laci.}{shift-buka}

\subsection{Selama shift}
Halaman menampilkan ringkasan yang bertambah setiap kali Anda menjual:

\begin{tabularx}{\textwidth}{@{}L{4.2cm}Y@{}}
  \toprule
  \textbf{Angka} & \textbf{Artinya} \\
  \midrule
  \field{Uang awal}       & Uang laci saat shift dibuka. \\
  \field{Transaksi}       & Jumlah transaksi pada shift ini. \\
  \field{Total penjualan} & Seluruh nilai penjualan shift ini. \\
  \field{Tunai}           & Bagian penjualan yang dibayar tunai. \\
  \field{Non-tunai}       & Bagian penjualan yang dibayar non-tunai (QRIS, debit, dan lain-lain). \\
  \field{Kas seharusnya}  & Uang tunai yang seharusnya ada di laci. \\
  \bottomrule
\end{tabularx}

\begin{itembox}{rumus kas seharusnya}
\[
\text{Kas seharusnya} = \text{Uang awal} + \text{Total penjualan} - \text{Non-tunai}
\]
Hanya penjualan \textbf{tunai} yang menambah uang di laci. Pembayaran QRIS,
debit, atau transfer tidak masuk laci.
\end{itembox}

\subsection{Menutup shift}

\begin{enumerate}[leftmargin=1.4em,itemsep=0.35em]
  \item Hitung uang tunai fisik di laci.
  \item Isi \field{Uang fisik dihitung} dengan hasil hitungan tersebut.
  \item Perhatikan baris \field{Selisih}: hijau berarti pas (Rp0), merah berarti ada
        selisih antara uang fisik dan \emph{kas seharusnya}.
  \item Isi \field{Catatan} bila ada selisih, agar dapat ditelusuri.
  \item Klik \btn{Tutup shift}.
\end{enumerate}

\begin{warnbox}{selisih kas}
Selisih muncul dari kesalahan kembalian, transaksi yang tidak tercatat, atau uang
awal yang salah diisi. Bila selisih muncul, jangan langsung mengubah angka ---
tulis catatan dan laporkan agar bisa ditelusuri.
\end{warnbox}

\screenshot{Halaman shift saat berjalan: penjualan, tunai, dan kas seharusnya.}{shift-terbuka}

\subsection{Riwayat shift}
Tabel \textbf{Riwayat shift} menampilkan kolom \field{Kasir}, \field{Buka},
\field{Penjualan}, \field{Seharusnya}, \field{Fisik}, dan \field{Selisih}. Shift
yang masih berjalan diberi label \textbf{terbuka}.

\begin{tipbox}{satu user, satu shift}
Satu pengguna hanya boleh memiliki satu shift terbuka. Untuk berganti kasir,
tutup shift terlebih dahulu, lalu kasir berikutnya membuka shift baru.
\end{tipbox}

% =====================================================================
\section{Kasir (POS)}
\label{sec:kasir}
\textbf{Judul halaman:} \emph{Kasir --- ``Transaksi penjualan \& penebusan resep''}

Ini halaman yang paling sering dipakai. Bila Anda belum membuka shift, halaman
ini akan menampilkan kartu \textbf{``Buka shift dulu''} --- buka shift terlebih
dahulu (Bagian~\ref{sec:shift}).

\subsection{Langkah melayani satu pembeli}

\begin{enumerate}[leftmargin=1.5em,itemsep=0.4em]
  \item \textbf{Cari produk.} Ketik nama pada kotak
        \field{``Scan barcode atau cari nama obat...''}. Bila ada barcode,
        arahkan pemindai ke kotak tersebut. Hasil pencarian menampilkan
        \field{Stok} dan harga per satuan.
  \item \textbf{Tambahkan ke keranjang.} Klik salah satu tombol satuan:
        \btn{+ tablet} untuk satuan dasar, atau \btn{+ strip (Rp6.500)} untuk
        satuan besar. Produk yang wajib resep diberi label \textbf{Resep}.
  \item \textbf{Atur jumlah.} Pada keranjang, gunakan tombol \textbf{+} dan
        \textbf{--}, atau ikon tempat sampah untuk menghapus satu baris.
  \item \textbf{Pilih pelanggan} pada kotak \field{Pelanggan} (bawaan:
        ``--- Umum ---'').
  \item \textbf{Isi diskon/pajak} bila ada pada kolom \field{Diskon} dan
        \field{Pajak}.
  \item \textbf{Pilih metode pembayaran} pada kotak \field{Metode}:
        \code{cash}, \code{debit}, \code{credit}, \code{qris},
        \code{ewallet}, atau \code{transfer}.
  \item \textbf{Isi jumlah uang yang diterima} pada kotak \field{Dibayar}.
        Untuk tunai, sistem otomatis menampilkan \field{Kembalian}.
  \item \textbf{Klik} \btn{Bayar Rp...}. Tunggu sampai muncul kartu hijau
        ``Transaksi \dots\ berhasil''.
  \item \textbf{Klik} \btn{Cetak struk} untuk mencetak, lalu \btn{Tutup} untuk
        melayani pembeli berikutnya.
\end{enumerate}

\begin{tabularx}{\textwidth}{@{}L{3.4cm}Y@{}}
  \toprule
  \textbf{Ringkasan nota} & \textbf{Keterangan} \\
  \midrule
  \field{Subtotal}  & Jumlah sebelum diskon dan pajak. \\
  \field{Total}     & Jumlah akhir yang harus dibayar. \\
  \field{Kembalian} & Uang kembali (ditampilkan untuk pembayaran tunai). \\
  \bottomrule
\end{tabularx}

\begin{tipbox}{memakai satuan besar}
Menjual 1 strip dari produk yang satuannya sudah diatur akan otomatis mengurangi
stok 10 tablet dan memakai harga strip. Tidak perlu menghitung tablet satu per
satu.
\end{tipbox}

\screenshot{Kasir: keranjang, total, dan kembalian sebelum pembayaran.}{kasir-keranjang}

\screenshot{Konfirmasi transaksi berhasil dengan nomor nota dan tombol \btn{Cetak struk}.}{kasir-berhasil}

\subsection{Struk}
Struk dicetak berukuran 80~mm (printer termal). Isinya: nama dan alamat apotek,
nomor nota, tanggal, nama kasir, nama pelanggan (bila ada), rincian item,
subtotal, diskon, pajak, \textbf{TOTAL}, metode pembayaran, kembalian (bila
tunai), catatan kaki, dan nama apoteker penanggung jawab.

\begin{warnbox}{sebelum menutup transaksi}
Periksa kembali item dan jumlah pada keranjang. Setelah nota tersimpan, transaksi
tidak dapat dihapus dari antarmuka. Bila terjadi kesalahan, catat nomor nota
(\code{INV-...}) dan laporkan ke tim IT.
\end{warnbox}

% =====================================================================
\section{Transaksi}
\label{sec:transaksi}
\textbf{Judul halaman:} \emph{Transaksi --- ``Riwayat penjualan dan cetak ulang struk''}

\subsection{Menyaring riwayat}

\begin{tabularx}{\textwidth}{@{}L{4cm}Y@{}}
  \toprule
  \textbf{Kontrol} & \textbf{Keterangan} \\
  \midrule
  \field{Sejak tanggal} & Menampilkan transaksi mulai tanggal tertentu (bawaan: hari ini). \\
  \field{Semua transaksi} & Centang untuk menampilkan seluruh riwayat tanpa batas tanggal. \\
  \bottomrule
\end{tabularx}

Di atas tabel tertera ringkasan \field{Transaksi} (jumlah transaksi) dan
\field{Omzet} (total nilai penjualan).

\screenshot{Riwayat transaksi dengan filter tanggal dan total omzet.}{transaksi}

\subsection{Mencetak ulang struk}
\begin{enumerate}[leftmargin=1.4em,itemsep=0.35em]
  \item Temukan nota pada tabel (kolom \field{Kode}, \field{Waktu},
        \field{Kasir}, \field{Item}, dan \field{Total}).
  \item Klik \btn{Struk} pada baris tersebut. Pratinjau muncul di panel kanan.
  \item Klik \btn{Cetak struk}.
\end{enumerate}

Nota yang dibatalkan ditandai label merah \textbf{void} pada kolom Kode.

% =====================================================================
\screenshot{Pratinjau struk 80~mm dan tombol \btn{Cetak struk}.}{transaksi-struk}

\section{Resep}
\label{sec:resep}
\textbf{Judul halaman:} \emph{Resep --- ``Pencatatan resep dokter dan penebusan obat''}

\screenshot{Daftar resep beserta statusnya.}{resep}

\subsection{Mencatat resep}

\begin{enumerate}[leftmargin=1.4em,itemsep=0.35em]
  \item Klik \btn{Input resep}.
  \item Isi \field{Nama pasien} (wajib).
  \item Pilih \field{Dokter} (bila ada) dan \field{Pelanggan} (bila terdaftar).
  \item Pada blok \field{Item resep}, klik \btn{Obat} lalu pilih obat, isi
        \field{Jumlah}, dan tulis \field{Signa} (mis. ``3 x 1 sehari'').
  \item Klik \btn{Simpan resep}.
\end{enumerate}

\begin{tabularx}{\textwidth}{@{}L{2.8cm}Y@{}}
  \toprule
  \textbf{Status resep} & \textbf{Artinya} \\
  \midrule
  \textbf{pending}   & Baru dicatat, belum diserahkan. \\
  \textbf{dispensed} & Obat sudah diserahkan ke pasien. \\
  \textbf{cancelled} & Resep dibatalkan. \\
  \bottomrule
\end{tabularx}

\screenshot{Formulir input resep: pasien, dokter, dan item resep.}{resep-form}

\begin{tipbox}{resep lalu penjualan}
Untuk menebus resep, catat resepnya di halaman ini, lalu lakukan penjualan di
halaman Kasir seperti biasa. Produk berlabel \textbf{Resep} adalah obat yang
memerlukan resep.
\end{tipbox}

\subsection{Catatan}
Mengubah status resep (\emph{pending/ dispensed/ cancelled}) belum tersedia
lewat tombol di halaman ini; saat ini melalui tim IT/API.

% =====================================================================
\section{Pengaturan}
\textbf{Judul halaman:} \emph{Pengaturan --- ``Profil apotek, apoteker penanggung jawab, dan struk''}

Nilai di halaman ini dipakai pada struk dan identitas apotek. Hanya
\textbf{owner} dan \textbf{admin} yang dapat menyimpan perubahan.

\begin{tabularx}{\textwidth}{@{}L{4.4cm}L{1.5cm}Y@{}}
  \toprule
  \textbf{Kolom isian} & \textbf{Wajib} & \textbf{Keterangan} \\
  \midrule
  \field{Nama apotek} & Ya & Tampil di bagian atas struk. \\
  \field{Telepon} & Tidak & Tampil sebagai ``Telp. \dots'' pada struk. \\
  \field{Alamat}  & Tidak & Tampil di bawah nama apotek. \\
  \field{NPWP}    & Tidak & Identitas pajak apotek. \\
  \field{PPN (\%)} & Tidak & Persentase pajak bawaan. \\
  \field{Apoteker penanggung jawab} & Tidak & Nama yang muncul di kaki struk. \\
  \field{No. SIPA} & Tidak & Nomor izin apoteker penanggung jawab. \\
  \field{Catatan kaki struk} & Tidak & Ucapan penutup pada struk, mis. ``Terima kasih, semoga lekas sembuh!''. \\
  \bottomrule
\end{tabularx}

Klik \btn{Simpan}. Bila berhasil, muncul pesan ``Pengaturan tersimpan.''

\begin{tipbox}{cek hasilnya di struk}
Setelah menyimpan, cetak satu struk uji dari halaman Kasir untuk memastikan nama,
alamat, dan catatan kaki sudah benar.
\end{tipbox}

% =====================================================================
\screenshot{Halaman pengaturan: profil apotek dan catatan kaki struk.}{pengaturan}

\section{Log audit}
\label{sec:audit}

\app{} mencatat aktivitas penting secara otomatis ke berkas harian. Yang dicatat
antara lain: login berhasil dan gagal, transaksi penjualan, buka/tutup shift,
perubahan pengaturan, penyesuaian stok, dan perubahan data lain
(create/update/delete).

\begin{tabularx}{\textwidth}{@{}L{4.6cm}Y@{}}
  \toprule
  \textbf{Aspek} & \textbf{Keterangan} \\
  \midrule
  Siapa yang bisa melihat & Hanya \textbf{owner} dan \textbf{admin}. \\
  Bentuk penyimpanan & Berkas per hari bernama \code{audit-YYYY-MM-DD.log} pada server. \\
  Isi setiap baris & Waktu, pengguna, aksi, metode/alamat, hasil, dan ringkasan data terkait. \\
  Keamanan & Kata sandi dan token tidak pernah ditulis ke log. \\
  Masa simpan & Berkas lama dihapus otomatis (bawaan 90 hari). \\
  Akses & Melalui tim IT/API (\code{GET /api/audit}, \code{GET /api/audit/days}). \\
  \bottomrule
\end{tabularx}

\begin{itembox}{untuk apa ini}
Log audit membantu menelusuri kejadian, misalnya ``siapa yang mengubah harga
produk ini'' atau ``kapan shift ditutup''. Log tidak dapat diubah dari aplikasi.
\end{itembox}

% =====================================================================
\section{Peran dan hak akses}
\label{sec:rbac}

Tabel berikut merangkum siapa yang boleh melakukan apa. Hak akses ditegakkan di
server, bukan hanya disembunyikan di tampilan.

\begin{longtable}{@{}L{5.4cm}L{9.4cm}@{}}
  \toprule
  \textbf{Tindakan} & \textbf{Peran yang diizinkan} \\
  \midrule
  \endhead
  Melihat semua halaman (baca) & Semua peran yang sudah masuk \\
  Menambah/mengubah katalog produk & owner, admin, pharmacist \\
  Membuat PO, menerima barang, membatalkan PO & owner, admin, pharmacist \\
  Penyesuaian stok (opname/rusak/expired) & owner, admin, pharmacist \\
  Mencatat resep & owner, admin, pharmacist, cashier \\
  Melakukan penjualan & Semua peran \emph{dan} wajib punya shift terbuka \\
  Buka/tutup shift & Semua peran (tutup hanya oleh pemilik shift atau owner) \\
  Mendaftarkan pelanggan & Semua peran \\
  Mengubah kategori, produsen, supplier, dokter & owner, admin, pharmacist \\
  Mengubah pengaturan apotek & owner, admin \\
  Melihat log audit & owner, admin \\
  \bottomrule
\end{longtable}

% =====================================================================
\section{Pesan yang sering muncul dan solusinya}

\begin{longtable}{@{}L{6.0cm}L{8.8cm}@{}}
  \toprule
  \textbf{Pesan} & \textbf{Penyebab dan solusi} \\
  \midrule
  \endhead
  \textbf{``Buka shift kasir terlebih dahulu sebelum bertransaksi''}
  & Anda menekan \btn{Bayar} tanpa shift terbuka. Buka menu \textbf{Shift Kasir} lalu klik \btn{Buka shift}. \\
  \addlinespace
  \textbf{``Keranjang masih kosong''}
  & Belum ada item yang ditambahkan. Cari produk dan klik tombol satuan \btn{+ ...} lebih dulu. \\
  \addlinespace
  \textbf{``Stok tidak mencukupi untuk salah satu produk''}
  & Jumlah yang diminta melebihi stok yang ada. Periksa \textbf{Persediaan} $\rightarrow$ tab \textbf{Batch \& Expiry}; bila barang sudah datang, pastikan penerimaan barang sudah dilakukan di halaman Pembelian. \\
  \addlinespace
  \textbf{``Stok berubah karena transaksi lain. Silakan coba lagi.''}
  & Dua kasir menjual produk yang sama pada saat yang hampir bersamaan. Ulangi transaksi --- stok sudah diperbarui. \\
  \addlinespace
  \textbf{``Jumlah pembayaran kurang dari total''}
  & Nilai pada \field{Dibayar} lebih kecil dari \field{Total}. Perbaiki jumlahnya. \\
  \addlinespace
  \textbf{``Email atau password salah''}
  & Periksa ejaan email dan password. Perhatikan huruf besar/kecil. \\
  \addlinespace
  \textbf{``Terlalu banyak percobaan login. Coba lagi nanti.''}
  & Pengaman dari percobaan login berulang. Tunggu beberapa menit, lalu coba lagi. \\
  \addlinespace
  \textbf{``Tidak punya akses''}
  & Peran Anda tidak berhak untuk tindakan itu (lihat Bagian~\ref{sec:rbac}). Minta owner/admin/apoteker melakukannya. \\
  \addlinespace
  \textbf{``SKU sudah dipakai''}
  & Kode SKU harus unik. Gunakan kode lain, mis. tambahkan akhiran angka. \\
  \addlinespace
  \textbf{``Masih ada shift yang terbuka. Tutup dulu sebelum buka shift baru.''}
  & Satu pengguna hanya boleh punya satu shift terbuka. Tutup shift yang berjalan lebih dulu. \\
  \addlinespace
  \textbf{``Hanya pemilik shift yang boleh menutup shift ini''}
  & Anda mencoba menutup shift milik pengguna lain. Hanya pemilik shift (atau owner) yang boleh menutup. \\
  \addlinespace
  \textbf{``Pembelian sudah diterima''}
  & Barang pada PO tersebut sudah pernah diterima. Tidak perlu diterima dua kali. \\
  \bottomrule
\end{longtable}

% =====================================================================
\section{Praktik yang disarankan}

\begin{itemize}[leftmargin=1.4em,itemsep=0.45em]
  \item \textbf{Buka shift setiap mulai bekerja}, dan \textbf{tutup shift}
        sebelum pulang. Jangan meninggalkan shift terbuka.
  \item \textbf{Hitung uang fisik dengan tenang} saat menutup shift, dan tulis
        catatan bila ada selisih.
  \item \textbf{Isi nomor batch dan tanggal kedaluwarsa} setiap kali menerima
        barang, agar FEFO dan peringatan kedaluwarsa bekerja.
  \item \textbf{Periksa Dashboard dan tab Peringatan setiap pagi} untuk
        merencanakan pembelian.
  \item \textbf{Jual batch mendekati kedaluwarsa lebih dulu} --- sistem sudah
        mengurutkannya, tetapi pastikan stok fisik di rak mengikuti urutan yang
        sama.
  \item \textbf{Biasakan memilih pelanggan} pada transaksi, agar riwayat pembeli
        anggota tetap terekam.
  \item \textbf{Simpan/arsipkan struk} bila pembeli meminta salinan; cetak ulang
        tersedia di halaman Transaksi.
\end{itemize}

% =====================================================================
\appendix
\section{Data contoh (seed)}

Data berikut terpasang pada aplikasi contoh dan berguna untuk latihan.

\subsection{Apotek}
\begin{tabularx}{\textwidth}{@{}L{4.6cm}Y@{}}
  \toprule
  \textbf{Kolom} & \textbf{Nilai contoh} \\
  \midrule
  Nama apotek & Apoteku Sehat \\
  Alamat      & Jl. Melati No. 12, Jakarta \\
  Telepon     & 021-1234567 \\
  Apoteker    & Apt. Rina Sari \\
  No. SIPA    & SIPA-12345 \\
  Catatan kaki struk & Terima kasih, semoga lekas sembuh! \\
  \bottomrule
\end{tabularx}

\subsection{Produk}
\begin{tabularx}{\textwidth}{@{}L{2.4cm}L{4.6cm}L{2.6cm}Y@{}}
  \toprule
  \textbf{SKU} & \textbf{Nama} & \textbf{Golongan} & \textbf{Harga jual} \\
  \midrule
  OBT-0001 & Paracetamol 500mg & obat bebas & Rp700/tablet; Rp6.500/strip; Rp60.000/box \\
  OBT-0002 & Amoxicillin 500mg & obat keras (wajib resep) & Rp1.500/kapsul; Rp14.000/strip \\
  OBT-0003 & Vitamin C 500mg & obat bebas & Rp900/tablet; Rp25.000/tube \\
  OBT-0004 & OBH Combi Batuk Flu 60ml & obat bebas terbatas & Rp18.000/botol \\
  ALK-0001 & Termometer Digital & alkes & Rp45.000/pcs \\
  \bottomrule
\end{tabularx}

\subsection{Batch contoh}
\begin{tabularx}{\textwidth}{@{}L{2.6cm}L{4.4cm}L{3.4cm}Y@{}}
  \toprule
  \textbf{Batch} & \textbf{Produk} & \textbf{Kedaluwarsa} & \textbf{Jumlah} \\
  \midrule
  PCM-A1 & Paracetamol 500mg & +400 hari & 500 \\
  PCM-A2 & Paracetamol 500mg & +60 hari  & 120 \\
  AMX-B1 & Amoxicillin 500mg & +300 hari & 200 \\
  VTC-C1 & Vitamin C 500mg   & +500 hari & 300 \\
  OBH-D1 & OBH Combi 60ml    & +25 hari  & 24 \\
  THM-E1 & Termometer Digital & +900 hari & 10 \\
  \bottomrule
\end{tabularx}

\begin{tipbox}{berlatih FEFO}
Jual Paracetamol 1 tablet di halaman Kasir, lalu periksa
\textbf{Persediaan} $\rightarrow$ tab \textbf{Batch \& Expiry}. Batch
\code{PCM-A2} (kedaluwarsa lebih dekat) akan berkurang lebih dulu.
\end{tipbox}

\subsection{Akun dan master data lain}
\begin{itemize}[leftmargin=1.4em,itemsep=0.3em]
  \item \textbf{Akun} (semua berpassword \code{password123}):
    \begin{itemize}[leftmargin=1.2em,itemsep=0.15em,topsep=0.2em]
      \item \code{owner@apoteku.test} --- Budi Owner (owner)
      \item \code{kasir@apoteku.test} --- Sari Kasir (kasir)
      \item \code{apoteker@apoteku.test} --- Apt. Rina (apoteker)
    \end{itemize}
  \item \textbf{Kategori}: Analgesik, Antibiotik, Vitamin \& Suplemen, Batuk \&
        Flu, Alat Kesehatan.
  \item \textbf{Supplier}: PT Kimia Farma Trading, PT Enseval Putera.
  \item \textbf{Pelanggan}: Umum, Ibu Ani (anggota), Bapak Joko (anggota).
  \item \textbf{Dokter}: dr. Andi Wijaya, dr. Maya Putri, Sp.A.
\end{itemize}

\section{Daftar istilah}

\begin{tabularx}{\textwidth}{@{}L{4.2cm}Y@{}}
  \toprule
  \textbf{Istilah} & \textbf{Arti} \\
  \midrule
  \textbf{Batch} & Satu kelompok barang dari satu penerimaan, dengan nomor batch dan tanggal kedaluwarsa sendiri. \\
  \textbf{FEFO} & \emph{First Expired First Out} --- batch yang paling cepat kedaluwarsa dijual lebih dulu. \\
  \textbf{PO} & \emph{Purchase Order} --- pesanan pembelian ke supplier. \\
  \textbf{Satuan dasar} & Satuan terkecil untuk perhitungan stok (mis. tablet). \\
  \textbf{Konversi} & Jumlah satuan dasar di dalam satu satuan besar (mis. 1 strip = 10 tablet). \\
  \textbf{Shift} & Sesi kerja kasir tempat uang awal dan uang akhir laci dicatat. \\
  \textbf{Kas seharusnya} & Uang tunai yang seharusnya ada di laci pada akhir shift. \\
  \textbf{Selisih} & Perbedaan antara uang fisik yang dihitung dan kas seharusnya. \\
  \textbf{SKU} & Kode unik internal produk. \\
  \textbf{Struk} & Bukti transaksi berukuran 80 mm. \\
  \bottomrule
\end{tabularx}

\section{Menjalankan aplikasi (untuk tim IT)}

\begin{itembox}{cara cepat dengan Docker}
\begin{enumerate}[leftmargin=1.4em,itemsep=0.25em]
  \item Siapkan Docker, lalu dari folder proyek jalankan:
        \code{bun run docker:e2e} untuk menyalakan seluruh sistem (basis data,
        API, dan web) sekaligus menjalankan uji otomatis.
  \item Buka \code{http://localhost:8080} pada peramban.
  \item Untuk mematikan dan membersihkan: \code{bun run docker:e2e:down}.
\end{enumerate}
\end{itembox}

\begin{tabularx}{\textwidth}{@{}L{4.6cm}Y@{}}
  \toprule
  \textbf{Kebutuhan} & \textbf{Keterangan} \\
  \midrule
  Basis data & MySQL 8 atau MariaDB. \\
  Alamat bawaan & Web: \code{localhost:8080} (Docker) atau \code{localhost:5173} (pengembangan); API: \code{localhost:4000}. \\
  Migrasi \& data contoh & Dijalankan otomatis saat container pertama kali menyala. \\
  Pengujian & \code{bun run test:security}, \code{bun run test:e2e}, \code{bun run test:audit}, \code{bun run test:load}. \\
  \bottomrule
\end{tabularx}

\vfill
\begin{center}
  \small\textcolor{ink}{--- Akhir dokumen ---}
\end{center}

\end{document}
